% Chapter converted from Markdown
% Auto-generated - do not edit manually

\section*{Chapter 14 - Idea to Reality Engine (MIGE)}


\subsection*{Purpose}

This chapter documents the Memory-to-Idea Growth Engine (MIGE), the system that turns captured ideas into deployed systems. MIGE solves the fundamental problem introduced in Chapter 1: ideas die—there's no path from idea to reality, and execution is ad-hoc.

MIGE provides:
\item \textbf{Pipeline from spark to runtime} linking CMC memories, APOE chains, and deployment tooling
\item \textbf{Idea scoring} prioritizing ideas by impact, confidence, effort, and readiness
\item \textbf{Templates and assets} enabling rapid instantiation of common patterns
\item \textbf{Quality gates} ensuring quartet parity and validation at every stage

This chapter demonstrates that MIGE is not just a build system—it is the idea-to-reality engine that enables AIM-OS to evolve systematically. Without it, ideas remain unrealized, execution is ad-hoc, and quality is inconsistent.

\subsection*{Executive Summary}

MIGE transforms ideas into reality through an eight-stage pipeline: capture, classify, design, plan, build, validate, deploy, and learn. Ideas are scored by impact, confidence, effort, and readiness. Templates enable rapid instantiation. Quality gates ensure quartet parity. BTSM integration provides system context. HVCA enables three-mind coordination.

\textbf{Key Insight:} MIGE enables the "idea-to-reality" principle from Chapter 1. Without it, ideas remain unrealized and execution is ad-hoc. With it, every idea has a clear path to reality with quality validation at every stage.

\subsection*{Pipeline Overview}

MIGE operates through an eight-stage pipeline:

\subsubsection*{1. Capture}

\textbf{Process:} Ideas enter via Chat AI, CAS insights, or SIS retrospectives; stored as CMC atoms tagged \texttt{{type:'idea'}}

\textbf{Sources:}
\item \textbf{Chat AI:} User suggestions, feature requests
\item \textbf{CAS insights:} Anomaly-driven improvements
\item \textbf{SIS retrospectives:} Learning-driven enhancements

\textbf{Storage:} All ideas stored in CMC with tags for retrieval

\subsubsection*{2. Classify}

\textbf{Process:} Intent classification and CCS foreground decide category (feature, fix, experiment, doc)

\textbf{Categories:}
\item \textbf{Feature:} New functionality
\item \textbf{Fix:} Bug fixes or improvements
\item \textbf{Experiment:} Research or exploration
\item \textbf{Doc:} Documentation updates

\textbf{Mechanism:} CCS foreground analyzes intent → Classifies category → Routes to appropriate pipeline

\subsubsection*{3. Design}

\textbf{Process:} HHNI retrieves precedent; VIF sets confidence gate; SEG lists required anchors

\textbf{Steps:}
\item HHNI retrieves similar precedents
\item VIF sets confidence gate (typically ≥ 0.70)
\item SEG lists required evidence anchors

\textbf{Output:} Design specification with precedents, confidence, and evidence requirements

\subsubsection*{4. Plan}

\textbf{Process:} APOE builds orchestration chain; includes quality hooks (SDF-CVF) and evidence capture

\textbf{Components:}
\item APOE orchestration chain
\item Quality hooks (SDF-CVF checkpoints)
\item Evidence capture requirements

\textbf{Output:} Execution plan with steps, quality gates, and evidence requirements

\subsubsection*{5. Build}

\textbf{Process:} Code/templates generated; tests/examples implemented; tags updated

\textbf{Activities:}
\item Generate code from templates
\item Implement tests and examples
\item Update NL tags for quartet parity

\textbf{Output:} Built system with code, tests, docs, and tags

\subsubsection*{6. Validate}

\textbf{Process:} SDF-CVF suite ensures quartet parity; VIF recalculated

\textbf{Validation:}
\item SDF-CVF quartet parity check (P ≥ 0.90)
\item VIF confidence recalculation
\item Quality gate validation

\textbf{Output:} Validation results with confidence scores

\subsubsection*{7. Deploy}

\textbf{Process:} Application lifecycle tools push to staging/production; dashboards update

\textbf{Deployment:}
\item Push to staging environment
\item Run health checks
\item Deploy to production if validated
\item Update dashboards

\textbf{Output:} Deployed system with monitoring

\subsubsection*{8. Learn}

\textbf{Process:} SIS logs outcomes; CAS monitors impact; future ideas seeded

\textbf{Learning:}
\item SIS logs deployment outcomes
\item CAS monitors post-deployment impact
\item Future ideas seeded from learnings

\textbf{Output:} Learning artifacts and future ideas

This pipeline ensures systematic transformation from idea to reality with quality validation at every stage.

\subsection*{Runnable Examples (PowerShell)}

\subsubsection*{Example 1: Create Application Scaffold from Idea}
``\texttt{powershell
\section*{Create application scaffold from captured idea}
\_\_MATH\_BLOCK\_0\_\_true
        }
    } 
} | ConvertTo-Json -Depth 6

\_\_MATH\_BLOCK\_1\_\_app |
    Select-Object -ExpandProperty Content | ConvertFrom-Json

Write-Host "Application Created:"
Write-Host "  App ID: \_\_MATH\_BLOCK\_2\_\_result.app\_id)"
Write-Host "  Template: \_\_MATH\_BLOCK\_3\_\_result.template)"
Write-Host "  Quality Gates: \_\_MATH\_BLOCK\_4\_\_result.quality\_gates -join ', ')"
Write-Host "  Status: \_\_MATH\_BLOCK\_5\_\_result.status)"
``\texttt{

\subsubsection*{Example 2: Deploy Application to Staging}
}`\texttt{powershell
\section*{Deploy application to staging environment}
\_\_MATH\_BLOCK\_6\_\_true;
            monitoring=\_\_MATH\_BLOCK\_7\_\_true
        }
    } 
} | ConvertTo-Json -Depth 6

\_\_MATH\_BLOCK\_8\_\_deploy |
    Select-Object -ExpandProperty Content | ConvertFrom-Json

Write-Host "Deployment Status:"
Write-Host "  App ID: \_\_MATH\_BLOCK\_9\_\_result.app\_id)"
Write-Host "  Environment: \_\_MATH\_BLOCK\_10\_\_result.environment)"
Write-Host "  Health Checks: \_\_MATH\_BLOCK\_11\_\_result.health\_checks.status)"
Write-Host "  Deployment Status: \_\_MATH\_BLOCK\_12\_\_result.deployment\_status)"
``\texttt{

\subsubsection*{Example 3: Query Idea Pipeline Status}
}`\texttt{powershell
\section*{Query idea pipeline status and metrics}
\_\_MATH\_BLOCK\_13\_\_true
        }
    } 
} | ConvertTo-Json -Depth 6

\_\_MATH\_BLOCK\_14\_\_pipeline |
    Select-Object -ExpandProperty Content | ConvertFrom-Json

Write-Host "Idea Pipeline Status:"
Write-Host "  Total Ideas: \_\_MATH\_BLOCK\_15\_\_result.total\_ideas)"
Write-Host "  By Stage:"
\_\_MATH\_BLOCK\_16\_\_(\_\_MATH\_BLOCK\_17\_\_(\_\_MATH\_BLOCK\_18\_\_($result.avg\_time\_to\_deploy) days"
``\texttt{

\subsection*{Idea Scoring}

MIGE scores ideas using four dimensions:

\subsubsection*{Impact}

\textbf{Definition:} Expected value (users helped, risk reduced)

\textbf{Scoring:} High/medium/low impact based on:
\item Number of users affected
\item Risk reduction magnitude
\item Value delivered

\textbf{Use Case:} Prioritize high-impact ideas first

\subsubsection*{Confidence}

\textbf{Definition:} Derived from VIF, evidence coverage, precedent success

\textbf{Scoring:} High/medium/low confidence based on:
\item VIF confidence score (≥ 0.70 required)
\item Evidence coverage completeness
\item Precedent success rate

\textbf{Use Case:} Only proceed with high-confidence ideas

\subsubsection*{Effort}

\textbf{Definition:} Time/cost estimates from APOE chain

\textbf{Scoring:} Easy/medium/hard effort based on:
\item APOE chain complexity
\item Resource requirements
\item Time estimates

\textbf{Use Case:} Balance effort with impact

\subsubsection*{Readiness}

\textbf{Definition:} Availability of templates, tests, or existing components

\textbf{Scoring:} High/medium/low readiness based on:
\item Template availability
\item Test suite completeness
\item Component reusability

\textbf{Use Case:} Prioritize ready ideas for faster execution

\textbf{Composite Score:} }score = (0.4 × impact) + (0.3 × confidence) + (0.2 × effort) + (0.1 × readiness)\texttt{

Scores feed prioritization dashboards; only ideas above composite threshold advance.

\subsection*{Templates \& Assets}

MIGE uses templates and assets to accelerate development:

\subsubsection*{Template Library}

\textbf{Location:} }templates/mige/*\texttt{ contains scaffolds for:
\item Chat flows
\item MCP tools
\item Documentation
\item Dashboards

\textbf{Structure:} Every template includes quality gates: tests, docs stub, tag list

\subsubsection*{Template Instantiation}

\textbf{Process:} When instantiated, APOE ensures assets stored in CMC with proper tags and SEG anchors

\textbf{Steps:}
1. Select template from library
2. Instantiate with idea parameters
3. APOE validates quality gates
4. Store in CMC with tags
5. Create SEG anchors

\textbf{Output:} Instantiated system with quality gates satisfied

\textbf{Key Insight:} Templates enable rapid development while maintaining quality standards.

\subsection*{Integration Points}

MIGE integrates deeply with all AIM-OS systems:

\subsubsection*{CCS (Chapter 13)}

\textbf{CCS provides:} Background organizer provisions context; meta layer monitors quality  
\textbf{MIGE provides:} Ideas requiring coordination  
\textbf{Integration:} CCS ensures context available; meta monitors quality throughout pipeline

\textbf{Key Insight:} CCS coordinates MIGE execution. MIGE leverages CCS coordination.

\subsubsection*{CAS (Chapter 11)}

\textbf{CAS provides:} Anomaly monitoring post-deployment  
\textbf{MIGE provides:} Deployed systems requiring monitoring  
\textbf{Integration:} CAS monitors deployed systems; feeds SIS if improvements needed

\textbf{Key Insight:} CAS monitors MIGE outputs. MIGE benefits from CAS awareness.

\subsubsection*{SIS (Chapter 12)}

\textbf{SIS provides:} Deployment retrospectives and template refinement  
\textbf{MIGE provides:} Deployment outcomes for learning  
\textbf{Integration:} SIS consumes deployment retrospectives to refine templates

\textbf{Key Insight:} SIS improves MIGE templates. MIGE provides learning data to SIS.

\subsubsection*{SDF-CVF (Chapter 10)}

\textbf{SDF-CVF provides:} Quartet parity enforcement  
\textbf{MIGE provides:} Systems requiring quality validation  
\textbf{Integration:} SDF-CVF enforces quartet parity before release

\textbf{Key Insight:} SDF-CVF validates MIGE outputs. MIGE ensures quality through SDF-CVF.

\subsubsection*{ARD (Chapter 15)}

\textbf{ARD provides:} Research findings requiring execution  
\textbf{MIGE provides:} Execution pipeline for research findings  
\textbf{Integration:} ARD uses MIGE outputs as execution targets for deeper research findings

\textbf{Key Insight:} ARD generates research. MIGE executes research findings.

\textbf{Overall Insight:} MIGE is not isolated—it integrates with all systems to enable idea-to-reality transformation. Every system benefits from systematic execution.

\subsection*{Failure Modes \& Safeguards}

MIGE handles multiple failure scenarios:

\subsubsection*{Idea Backlog Overload}

\textbf{Scenario:} Too many ideas overwhelm the pipeline

\textbf{Safeguard:} Throttle intake; cluster ideas; auto-close duplicates

\textbf{Process:}
1. Detect backlog threshold exceeded
2. Throttle new idea intake
3. Cluster similar ideas
4. Auto-close duplicates

\textbf{Prevention:} Capacity limits, clustering algorithms, duplicate detection

\subsubsection*{Template Mismatch}

\textbf{Scenario:} Template doesn't fit idea requirements

\textbf{Safeguard:} Escalate to design review; create new template; update library

\textbf{Process:}
1. Detect template mismatch
2. Escalate to design review
3. Create new template if needed
4. Update template library

\textbf{Prevention:} Template validation, design review process

\subsubsection*{Deployment Failure}

\textbf{Scenario:} Deployment fails in production

\textbf{Safeguard:} Automatic rollback; capture logs; open remediation task via SIS

\textbf{Process:}
1. Detect deployment failure
2. Automatic rollback to previous version
3. Capture failure logs
4. Open remediation task via SIS

\textbf{Prevention:} Health checks, rollback procedures, monitoring

\subsubsection*{Quality Regression}

\textbf{Scenario:} Quality degrades after deployment

\textbf{Safeguard:} Block release; rerun improvements; update VIF; record in SEG

\textbf{Process:}
1. Detect quality regression
2. Block release
3. Rerun improvements
4. Update VIF confidence
5. Record in SEG

\textbf{Prevention:} Quality gates, regression testing, monitoring

Each failure mode has documented safeguards that preserve quality and enable recovery.

\subsection*{Ops Runbook}
1. Review idea queue; ensure metadata complete.
2. Approve ideas with high impact/confidence; assign owners.
3. Trigger APOE chain to generate plan + tasks.
4. Track progress via MIGE dashboard (build status, validation, deployment).
5. After deployment, run post-release checklist (quality metrics, user feedback, evidence logging).

\subsection*{Continuous Learning}
\item Post-mortems update scoring weights and templates.
\item Successful deployments spawn "pattern cards" stored in HHNI for faster future execution.
\item Failed experiments captured in SEG to warn future planners.
\item Metrics feed into SIS to propose improvements (better templates, gating logic).

\subsection*{Bitemporal Total System Map (BTSM) Integration}
MIGE leverages the BTSM to understand system context:

\item \textbf{System Inventory:} BTSM provides living inventory of every subsystem, dependency, and policy pack. Each node carries Minimal-Perfect-Details (MPD) including purpose, capabilities, interfaces, dependencies, and lifecycle.

\item \textbf{Blast Radius Analysis:} Before implementing ideas, MIGE queries BTSM to understand impact radius. Changes are validated against affected systems to prevent unintended consequences.

\item \textbf{Dependency Resolution:} BTSM dependency graphs enable MIGE to resolve dependencies automatically. Ideas that require unavailable dependencies are flagged or deferred.

\item \textbf{Temporal Replay:} BTSM's bitemporal edges enable replaying system state at any moment. MIGE uses this for debugging and understanding historical context.

\subsection*{Harmonised Verifiable Cognitive Architecture (HVCA)}
MIGE employs HVCA's three-mind neuro-symbolic loop:

\item \textbf{Mind 1 (Meta-Optimizer):} Shapes the vision tensor from human seed ideas. Optimizes idea formulation for maximum impact and feasibility.

\item \textbf{Mind 2 (Context Retriever):} Gathers context slices using DVNS and REX-RAG. Retrieves relevant precedents, patterns, and knowledge from HHNI.

\item \textbf{Mind 3 (Constraint Enforcer):} Ensures feasibility using symbolic reasoning and MCCA scores. Validates ideas against system constraints, budgets, and policies.

\item \textbf{Coordination:} All three minds coordinate through APOE, with every exchange emitting VIF evidence for auditability.

\subsection*{Quality Gates \& Validation}
MIGE enforces quality at every pipeline stage:

\item \textbf{Vision Gate:} }g\_vision\_fit\texttt{ (>= 0.90) ensures ideas align with system vision and goals.

\item \textbf{Trunk Gate:} }g\_trunk\_coherence\texttt{ and }g\_scope\_coverage\texttt{ validate ideas fit within system architecture.

\item \textbf{Variant Gate:} }g\_variant\_parity\texttt{ ensures design variants maintain consistency.

\item \textbf{Budget Gate:} }g\_budget\_guard\texttt{ prevents resource overruns.

\item \textbf{Quartet Parity:} SDF-CVF ensures code, docs, tests, and traces maintain parity (P >= 0.90).

\subsection*{Connection to Other Chapters}

MIGE connects to all AIM-OS systems:

\item \textbf{Chapter 1 (The Great Limitation):} MIGE addresses "ideas die" by enabling systematic execution
\item \textbf{Chapter 2 (The Vision):} MIGE enables the "idea-to-reality" principle from the universal interface
\item \textbf{Chapter 3 (The Proof):} MIGE validates execution through quality gates
\item \textbf{Chapter 5 (CMC):} MIGE stores all ideas and artifacts in CMC for durability
\item \textbf{Chapter 6 (HHNI):} MIGE uses HHNI for precedent retrieval
\item \textbf{Chapter 7 (VIF):} MIGE uses VIF for confidence gating
\item \textbf{Chapter 8 (APOE):} MIGE uses APOE for orchestration chains
\item \textbf{Chapter 9 (SEG):} MIGE uses SEG for evidence anchors
\item \textbf{Chapter 10 (SDF-CVF):} MIGE uses SDF-CVF for quality validation
\item \textbf{Chapter 11 (CAS):} MIGE uses CAS for post-deployment monitoring
\item \textbf{Chapter 12 (SIS):} MIGE uses SIS for learning and template refinement
\item \textbf{Chapter 13 (CCS):} MIGE uses CCS for coordination
\item \textbf{Chapter 15 (ARD):} MIGE executes ARD research findings

\textbf{Key Insight:} MIGE is the execution engine that transforms ideas into reality. Without MIGE, ideas remain unrealized and execution is ad-hoc.

\subsection*{MIGE Architecture \& System Design}

MIGE implements a comprehensive framework for transforming ideas into deployed systems through systematic pipeline execution.

\subsubsection*{Core Pipeline Architecture}

\textbf{Eight-Stage Pipeline:}
1. \textbf{Capture:} Ideas enter via multiple sources, stored as CMC atoms
2. \textbf{Classify:} Intent classification routes ideas to appropriate pipelines
3. \textbf{Design:} Precedent retrieval, confidence gating, evidence requirements
4. \textbf{Plan:} APOE orchestration chains with quality hooks
5. \textbf{Build:} Code generation, test implementation, tag updates
6. \textbf{Validate:} SDF-CVF quartet parity, VIF recalculation
7. \textbf{Deploy:} Application lifecycle deployment with health checks
8. \textbf{Learn:} SIS retrospectives, CAS monitoring, future idea seeding

\textbf{Pipeline Characteristics:}
\item \textbf{Linear Flow:} Each stage must complete before next begins
\item \textbf{Quality Gates:} Validation at every stage prevents regressions
\item \textbf{Evidence Tracking:} All artifacts linked to SEG evidence anchors
\item \textbf{Confidence Tracking:} VIF confidence updated throughout pipeline
\item \textbf{Audit Trail:} Complete bitemporal tracking via CMC

\subsubsection*{Idea Scoring System}

\textbf{Four-Dimensional Scoring:}
\item \textbf{Impact (40\%):} Expected value, users helped, risk reduced
\item \textbf{Confidence (30\%):} VIF score, evidence coverage, precedent success
\item \textbf{Effort (20\%):} APOE chain complexity, resource requirements
\item \textbf{Readiness (10\%):} Template availability, test completeness

\textbf{Scoring Formula:} }score = (0.4 × impact) + (0.3 × confidence) + (0.2 × effort) + (0.1 × readiness)\texttt{

\textbf{Thresholds:}
\item \textbf{High Priority:} score ≥ 0.80
\item \textbf{Medium Priority:} 0.60 ≤ score < 0.80
\item \textbf{Low Priority:} score < 0.60

\subsubsection*{Template System Architecture}

\textbf{Template Library Structure:}
\item \textbf{Location:} }templates/mige/*\texttt{ organized by category
\item \textbf{Categories:} Chat flows, MCP tools, Documentation, Dashboards, Services
\item \textbf{Structure:} Each template includes code, tests, docs, tags, quality gates

\textbf{Template Instantiation Process:}
1. Select template from library
2. Instantiate with idea parameters
3. APOE validates quality gates
4. Store in CMC with tags
5. Create SEG anchors
6. Update HHNI nodes

\textbf{Template Evolution:}
\item Successful deployments spawn pattern cards
\item Pattern cards stored in HHNI for faster future execution
\item Templates refined based on SIS retrospectives
\item Failed experiments captured in SEG to warn future planners

\subsubsection*{Quality Gate Architecture}

\textbf{Gate Types:}
\item \textbf{Vision Gate:} }g\_vision\_fit\texttt{ (>= 0.90) - Idea aligns with system vision
\item \textbf{Trunk Gate:} }g\_trunk\_coherence\texttt{, }g\_scope\_coverage\texttt{ - Fits architecture
\item \textbf{Variant Gate:} }g\_variant\_parity\texttt{ - Design variants maintain consistency
\item \textbf{Budget Gate:} }g\_budget\_guard` - Prevents resource overruns
\item \textbf{Quartet Parity:} SDF-CVF ensures code/docs/tests/traces (P >= 0.90)

\textbf{Gate Enforcement:}
\item Gates checked at every pipeline stage
\item Failure blocks progression until resolved
\item Gate results recorded in VIF witnesses
\item Gate violations trigger SIS improvement dreams

\subsubsection*{BTSM Integration Architecture}

\textbf{System Inventory:}
\item BTSM provides living inventory of every subsystem
\item Each node carries Minimal-Perfect-Details (MPD)
\item MPD includes purpose, capabilities, interfaces, dependencies, lifecycle

\textbf{Blast Radius Analysis:}
\item MIGE queries BTSM before implementation
\item Impact radius calculated for all affected systems
\item Changes validated against affected systems
\item Prevents unintended consequences

\textbf{Dependency Resolution:}
\item BTSM dependency graphs enable automatic resolution
\item Ideas requiring unavailable dependencies flagged
\item Dependency availability checked before planning
\item Deferred ideas tracked for future execution

\subsubsection*{HVCA Integration Architecture}

\textbf{Three-Mind Neuro-Symbolic Loop:}
\item \textbf{Mind 1 (Meta-Optimizer):} Shapes vision tensor from human seed ideas
\item \textbf{Mind 2 (Context Retriever):} Gathers context using DVNS and REX-RAG
\item \textbf{Mind 3 (Constraint Enforcer):} Ensures feasibility using symbolic reasoning

\textbf{Coordination:}
\item All three minds coordinate through APOE
\item Every exchange emits VIF evidence for auditability
\item MCCA scores validate constraint satisfaction
\item Vision tensor optimized for maximum impact and feasibility

\subsection*{Real-World MIGE Operations}

\subsubsection*{Case Study: Rapid Feature Development}

\textbf{Scenario:} Build new MCP tool integration feature in 3 days.

\textbf{MIGE Pipeline Execution:}
1. \textbf{Capture:} Feature idea captured via Chat AI, stored as CMC atom
2. \textbf{Classify:} Classified as "feature" category, routed to feature pipeline
3. \textbf{Design:} HHNI retrieved similar MCP tool precedents, VIF set confidence gate (0.85), SEG listed required anchors
4. \textbf{Plan:} APOE created orchestration chain with quality hooks, evidence capture requirements
5. \textbf{Build:} Code generated from MCP tool template, tests implemented, NL tags updated
6. \textbf{Validate:} SDF-CVF validated quartet parity (0.92), VIF recalculated confidence (0.88)
7. \textbf{Deploy:} Deployed to staging, health checks passed, deployed to production
8. \textbf{Learn:} SIS logged outcomes, CAS monitored impact, future ideas seeded

\textbf{Outcome:} Feature completed in 2.5 days with all quality gates passing, zero regressions, complete documentation.

\textbf{Metrics:}
\item \textbf{Development Time:} 2.5 days (target: 3 days) ✅
\item \textbf{Quality Gates:} All passing ✅
\item \textbf{Quartet Parity:} 0.92 (target: ≥0.90) ✅
\item \textbf{Regressions:} 0 (zero regressions) ✅
\item \textbf{Documentation:} Complete ✅

\textbf{Key Learnings:}
\item MIGE accelerates idea-to-deployment pipeline
\item Template system enables rapid development
\item Quality gates prevent regressions
\item BTSM integration prevents unintended consequences

\subsubsection*{Case Study: System Refactoring}

\textbf{Scenario:} Refactor legacy system to modern architecture.

\textbf{MIGE Pipeline Execution:}
1. \textbf{Capture:} Refactoring idea captured via CAS anomaly detection
2. \textbf{Classify:} Classified as "fix" category, routed to improvement pipeline
3. \textbf{Design:} BTSM analyzed blast radius, identified affected systems, VIF set confidence gate (0.75)
4. \textbf{Plan:} APOE created phased orchestration chain with rollback checkpoints
5. \textbf{Build:} Refactored code generated, comprehensive tests implemented, migration scripts created
6. \textbf{Validate:} SDF-CVF validated quartet parity (0.91), blast radius verified, rollback tested
7. \textbf{Deploy:} Phased deployment to staging, validation at each phase, production deployment
8. \textbf{Learn:} SIS logged outcomes, CAS monitored for regressions, pattern card created

\textbf{Outcome:} System refactored successfully with zero downtime, complete test coverage, documented migration path.

\textbf{Metrics:}
\item \textbf{Refactoring Time:} 2 weeks (target: 2 weeks) ✅
\item \textbf{Downtime:} 0 (zero downtime) ✅
\item \textbf{Test Coverage:} 100\% ✅
\item \textbf{Blast Radius:} All affected systems identified ✅
\item \textbf{Rollback Tested:} Verified ✅

\textbf{Key Learnings:}
\item BTSM enables accurate blast radius analysis
\item Phased deployment reduces risk
\item Rollback checkpoints enable safe refactoring
\item Pattern cards accelerate future refactoring

\subsection*{MIGE Performance Characteristics}

\subsubsection*{Pipeline Latency}

\textbf{Stage Timings:}
\item \textbf{Capture:} <1 second (immediate storage)
\item \textbf{Classify:} <5 seconds (intent analysis)
\item \textbf{Design:} <30 seconds (precedent retrieval, confidence calculation)
\item \textbf{Plan:} <2 minutes (APOE chain generation)
\item \textbf{Build:} Variable (depends on complexity, typically 5-30 minutes)
\item \textbf{Validate:} <5 minutes (SDF-CVF suite execution)
\item \textbf{Deploy:} <10 minutes (staging deployment, health checks)
\item \textbf{Learn:} <1 minute (SIS logging, CAS monitoring)

\textbf{Total Pipeline Time:} Typically 15-60 minutes for simple ideas, 2-8 hours for complex ideas.

\subsubsection*{Throughput Requirements}

\textbf{Idea Processing:}
\item \textbf{Capture Rate:} 100+ ideas/day
\item \textbf{Classification Rate:} 50+ ideas/hour
\item \textbf{Design Rate:} 20+ ideas/hour
\item \textbf{Planning Rate:} 10+ ideas/hour
\item \textbf{Build Rate:} 5+ ideas/hour
\item \textbf{Validation Rate:} 10+ validations/hour
\item \textbf{Deployment Rate:} 5+ deployments/hour

\textbf{Key Insight:} MIGE throughput requirements enable high-volume idea processing while maintaining quality.

\subsubsection*{Quality Metrics}

\textbf{Success Rates:}
\item \textbf{Pipeline Completion:} 85\%+ ideas complete pipeline successfully
\item \textbf{Quality Gate Pass:} 90\%+ ideas pass all quality gates
\item \textbf{Deployment Success:} 95\%+ deployments successful
\item \textbf{Zero Regression:} 98\%+ deployments introduce zero regressions

\textbf{Key Insight:} MIGE quality metrics ensure high success rates while maintaining quality standards.
